\documentclass[11pt]{article}

\usepackage[margin=1in]{geometry}
\usepackage[T1]{fontenc}
\usepackage[utf8]{inputenc}
\usepackage{lmodern}
\usepackage{microtype}
\usepackage{amsmath,amssymb,amsthm,amscd,mathtools}
\usepackage{booktabs}
\usepackage{enumitem}
\usepackage[hidelinks]{hyperref}
\usepackage[nameinlink,capitalise]{cleveref}
\hypersetup{
  pdftitle={Modal Triplet Theory and the Hull-Strominger System: A Conditional Fixed-Point Correspondence and the q=79 Completion Boundary},
  pdfauthor={Peter Nero}
}

\newtheorem{theorem}{Theorem}[section]
\newtheorem{proposition}[theorem]{Proposition}
\newtheorem{lemma}[theorem]{Lemma}
\newtheorem{corollary}[theorem]{Corollary}
\theoremstyle{definition}
\newtheorem{definition}[theorem]{Definition}
\newtheorem{assumption}[theorem]{Assumption}
\newtheorem{example}[theorem]{Example}
\theoremstyle{remark}
\newtheorem{remark}[theorem]{Remark}

\newcommand{\dd}{\mathrm d}
\newcommand{\ii}{\mathrm i}
\newcommand{\cA}{\mathcal A}
\newcommand{\cB}{\mathcal B}
\newcommand{\cC}{\mathcal C}
\newcommand{\cD}{\mathcal D}
\newcommand{\cF}{\mathcal F}
\newcommand{\cH}{\mathcal H}
\newcommand{\cR}{\mathcal R}
\newcommand{\cU}{\mathcal U}
\newcommand{\cV}{\mathcal V}
\newcommand{\cY}{\mathcal Y}
\newcommand{\HS}{Hull--Strominger}
\newcommand{\MTT}{Modal Triplet Theory}
\newcommand{\Tr}{\operatorname{Tr}}
\newcommand{\tr}{\operatorname{tr}}
\newcommand{\Ran}{\operatorname{Ran}}
\newcommand{\Fix}{\operatorname{Fix}}
\newcommand{\Lip}{\operatorname{Lip}}

\title{\vspace{-1em}
Modal Triplet Theory and the Hull--Strominger System:\\[0.25em]
\large A Conditional Fixed-Point Correspondence and the \(q=79\)
Completion Boundary}
\author{Peter Nero}
\date{Version 3, September 2026}

\begin{document}
\maketitle

\begin{abstract}
The Hull--Strominger system couples a conformally balanced Hermitian metric,
holomorphic gauge bundles, Hermitian--Yang--Mills connections, torsion, and the
differential Green--Schwarz identity on one compact complex threefold.  This
paper asks a precise question: when does a fixed point of a Modal Triplet
Theory evolution determine a solution of that system?  We define a typed
bridge from an upstairs MTT configuration space to Hull--Strominger data and
prove exact and residual fixed-point descent theorems.  Exact descent requires
the bridge to intertwine the selected MTT flow with a lower geometric flow,
such as the Anomaly flow together with the required bundle evolutions.
Approximate intertwining bounds the lower generator; a full equation bound
additionally requires control of the residual-to-generator map.  A contextual
analysis of the selected \(q=79\) constructions separates the established
hidden projective rank-nine bundle and existential Hermitian--Yang--Mills
connection from the unconstructed physical visible rank-three bundle and
common endpoint.  Full Bott--Chern injectivity supplies an anomaly potential
once the same-source curvature pair exists, but does not supply positivity
or integral differential flux data.  A transverse Fu--Yau shortcut is excluded
for the required nonzero visible third Chern class; a full nontransverse
residual compiler is available conditionally.  These results establish a
typed completion contract, not compactification selection or a physical
action.  The connection-preserving continuum intertwiner remains open.
\end{abstract}

\section*{Revision note: Version 3}
\addcontentsline{toc}{section}{Revision note: Version 3}
\begin{description}[leftmargin=2.4cm,style=nextline]
\item[Supersedes]
The current Version 2 manuscript.  Its released DOI and the earlier revision
note below remain unchanged; this revision is not a new public release.
\item[Reason]
The evidence section stopped at rank-two analysis and did not distinguish
later hidden projective HYM existence from the still-open common physical
endpoint.  Its torsion convention also identified a degree-one operator with
a degree-two operator.
\item[Resolution]
The convention is corrected to \(\dd\dd^c=\ii\partial\bar\partial\).
Thirty-seven frozen source assignments are reviewed in context: source
typing, full-residual transfer, the successful hidden construction and its
excluded predecessors, Bott--Chern reduction, the nontransverse compiler,
and the integral visible-source boundary.  A fresh-prime arithmetic
correction is retained without discarding unaffected finite results.
Scoped consumer summaries identify the current Cohesive, quantum-mechanics,
and flux-paper owners of the geometric and operator interfaces.
\item[Retained]
Conditional fixed-point descent, the rejection of the former selection
potential, finite \(q=79\) and rank-two analytic results, and the established
hidden rank-nine object and existential HYM theorem at their declared tiers.
\item[Open boundary]
The selected visible twisted Prym datum, common visible--hidden chamber,
positive nontransverse solution, integral differential Green--Schwarz
trivialization, and same-source continuum intertwiner remain to be supplied.
\end{description}

\section*{Revision note: Version 2}
\addcontentsline{toc}{section}{Revision note: Version 2}
\begin{description}[leftmargin=2.4cm,style=nextline]
\item[Supersedes]
Version 2 supersedes Version 1 and its claim that MTT had uniquely selected a
non-K\"ahler heterotic compactification.

\item[Reason]
The earlier argument used a ``twisted differential'' even though the anomaly
equation generally gives \(\dd H\ne0\), inferred a contraction without a
verified contraction constant, and treated an indefinite constrained
functional as a coercive strictly convex selection potential.  It also
assembled Fu--Yau, Iwasawa, bundle, and MTT ingredients that had not been
constructed on one common carrier.

\item[Resolution]
This version removes the invalid selection functional and replaces it with a
typed flow-intertwining contract.  It proves exactly what follows from an
intertwiner, gives the corresponding residual estimate, and separates lower
Hull--Strominger existence from MTT source selection.

\item[Retained content]
The Hull--Strominger equations, the use of conformally balanced geometry and
Hermitian--Yang--Mills data, fixed-point methods as a possible bridge, and
Fu--Yau geometry as the strongest established lower-space target are retained.

\item[Remaining boundary]
MTT must still construct one selected \(q=79\) physical visible--hidden
background and prove that its upper evolution descends to the required
geometric and bundle flows with the same connections, traces, and global flux
data.
\end{description}

\tableofcontents

\section{The question and the answer}
\label{sec:intro}

\subsection{Why a fixed-point bridge is attractive}

The \HS{} system is not one equation.  It is an assembly of complex geometry,
gauge theory, a differential anomaly equation, and global flux data.  MTT, in
turn, is organized around admissible carriers, projected sectors, and
fixed-point evolution.  It is therefore natural to ask whether an MTT fixed
point can be transported into a heterotic fixed point.

The attraction of this idea should not obscure the logical order.  A fixed
point in one space is not automatically a fixed point in another.  One needs
a map between the spaces, and that map must respect the evolutions.  This
paper supplies that missing mathematical sentence:
\[
\begin{CD}
\cU_{\mathrm{MTT}} @>{R_\tau}>> \cU_{\mathrm{MTT}}\\
@V{\mathfrak B}VV @VV{\mathfrak B}V\\
\cY_{\mathrm{HS}} @>{S_\tau}>> \cY_{\mathrm{HS}} .
\end{CD}
\]
Here \(R_\tau\) is a selected MTT stabilization evolution,
\(S_\tau\) is a declared lower geometric evolution, and
\(\mathfrak B\) is a typed bridge.  The commuting square is an assumption
until its rows are derived from one source.

\subsection{What is proved}

The paper owns three limited results.
\begin{enumerate}[leftmargin=1.7em]
\item Exact flow intertwining sends an MTT fixed point to a lower fixed point.
\item A bounded intertwining defect bounds the lower generator, and bounds
      all equations only under a declared norm comparison.
\item Local uniqueness or selection requires an invariant contraction basin
      and cannot be inferred from fixed-point correspondence alone.
\end{enumerate}
These are mathematical bridge theorems.  They do not prove that MTT already
provides the hypotheses.

\subsection{Dependency map}

The argument has five layers:
\[
\begin{split}
\text{one Hull--Strominger carrier}
&\longrightarrow \text{lower residuals and lower flow}\\
&\longrightarrow \text{typed MTT-to-lower bridge}\\
&\longrightarrow \text{intertwining certificate}\\
&\longrightarrow \text{fixed-point descent}\\
&\longrightarrow \text{physical completion tests}.
\end{split}
\]
Sections~\ref{sec:hs} and \ref{sec:flow} explain the first two layers.
Sections~\ref{sec:bridge}--\ref{sec:selection} prove the bridge results.
Sections~\ref{sec:audit}--\ref{sec:frontier} state what survives physically.

\section{One Hull--Strominger object}
\label{sec:hs}

\subsection{Fields and convention}

Fix a compact complex threefold \(X\) with nowhere-vanishing holomorphic
\((3,0)\)-form \(\Omega\).  Let \(\omega\) be a positive Hermitian form,
\(\Phi\) a dilaton, and \(V_{\mathrm{vis}},V_{\mathrm{hid}}\) holomorphic
Hermitian gauge objects with unitary connections
\(A_{\mathrm{vis}},A_{\mathrm{hid}}\).  Fix also a metric connection
\(\nabla\) on \(TX\).  Its choice is part of the data, not a notation that may
be changed between equations.
For a twisted or projective object, its cocycle, determinant reduction, and
trace-free curvature convention are part of the datum.  The adjoint curvature
is globally defined although the vector transition maps may be projective.
An arbitrary coherent sheaf is not an instanton: local freeness, or an
explicitly different singular analytic problem, must be specified.

We use
\[
\dd^c=\frac{\ii}{2}(\bar\partial-\partial),\qquad
\dd\dd^c=\ii\partial\bar\partial,\qquad
H=\dd^c\omega
\]
as a convention for the torsion equation.  Numerical factors in \(\dd^c\)
and trace normalizations vary in the literature; every comparison below
presupposes one convention fixed throughout.

In this convention the first-order equations include
\begin{align}
\dd\!\left(\|\Omega\|_\omega\omega^2\right)&=0,
\label{eq:balanced}\\
F_a^{0,2}=0,\qquad F_a\wedge\omega^2&=0,
\quad a\in\{\mathrm{vis},\mathrm{hid}\},
\label{eq:gauge-hym}\\
R_\nabla^{0,2}=0,\qquad R_\nabla\wedge\omega^2&=0,
\label{eq:tangent-instanton}\\
\dd H&=\frac{\alpha'}4\left(
\tr R_\nabla\wedge R_\nabla
-\tr F_{\mathrm{vis}}\wedge F_{\mathrm{vis}}
-\tr F_{\mathrm{hid}}\wedge F_{\mathrm{hid}}\right).
\label{eq:bianchi}
\end{align}
The tangent-instanton row is included when required by the chosen
first-order equations-of-motion convention
\cite{Strominger1986,Hull1986,DelaOssaSvanes2014}.

\subsection{Global data are a separate row}

Writing \(H=\dd B+\) Chern--Simons terms is local notation.  Globally, the
\(B\)-field is gerbe data and the Green--Schwarz condition is differential
cohomological.  A topological equality of second Chern classes is necessary
in common settings but does not identify the differential four-form
representatives in \cref{eq:bianchi}.  Likewise, a smooth vector bundle with
the desired Chern classes is not yet a holomorphic stable bundle with an HYM
connection.

\begin{definition}[Complete lower datum]
\label{def:lower}
A point \(y\in\cY_{\mathrm{HS}}\) is a tuple
\[
y=(X,\Omega,\omega,\Phi,V_{\mathrm{vis}},V_{\mathrm{hid}},
A_{\mathrm{vis}},A_{\mathrm{hid}},\nabla,H,\mathfrak g_H)
\]
in which all objects live on the same \(X\), use one trace convention, and
\(\mathfrak g_H\) denotes the global gerbe or differential-cohomological flux
datum.  Gauge-equivalent tuples represent the same point.
\end{definition}

\begin{definition}[Hull--Strominger residual]
\label{def:residual}
Fix the differential-topological class and its admissible patching first.
After gauge fixing and choosing Sobolev completions, let
\(\cR_{\mathrm{HS}}(y)\) be the vector consisting of the left-hand sides of
\cref{eq:balanced,eq:gauge-hym,eq:tangent-instanton,eq:bianchi}, together
with the continuous patching equations in this class.  Discrete obstruction
classes are admission conditions, not extra coordinates in a Sobolev vector
space.  The torsion and dilaton identifications are imposed in the domain
(in particular, \(e^{-2\Phi}\) is proportional to \(\|\Omega\|_\omega\)
with the chosen normalization).  Thus
\[
\cR_{\mathrm{HS}}(y)=0
\]
means that every declared lower row is satisfied on one common tuple.
\end{definition}

The residual formulation is intentionally unforgiving.  It prevents a metric
from one construction, a bundle from another, and a topological identity from
a third from being advertised as one solution.

\section{The appropriate lower flow}
\label{sec:flow}

\subsection{Anomaly flow}

There is an established geometric flow designed for this setting.  In a fixed
holomorphic and bundle ansatz, the Anomaly flow evolves the positive
\((2,2)\)-form \(\|\Omega\|_\omega\omega^2\) by an equation of the form
\begin{equation}
\partial_\tau\!\left(\|\Omega\|_\omega\omega^2\right)
=\ii\partial\bar\partial\omega
-\frac{\alpha'}4\left(
\tr R_\nabla\wedge R_\nabla
-\tr F_{\mathrm{vis}}\wedge F_{\mathrm{vis}}
-\tr F_{\mathrm{hid}}\wedge F_{\mathrm{hid}}\right).
\label{eq:anomaly-flow}
\end{equation}
The exact analytic formulation depends on the selected tangent connection and
on whether bundle metrics are fixed or evolved.  The flow preserves the
conformally balanced condition under its hypotheses, and its stationary
points solve the anomaly equation.  Short-time existence is known in the
original setting, and convergence is known in important Fu--Yau ansatz
classes \cite{PhongPicardZhang2017,PhongPicardZhang2018}.

The existence of \cref{eq:anomaly-flow} is important for MTT because it gives
a genuine lower repair dynamics.  It does not prove that an MTT flow equals
it.  Nor does the metric flow by itself construct the holomorphic bundles or
their HYM metrics.  Those rows must be fixed consistently or coupled to
appropriate bundle heat flows.

\subsection{A lower semigroup is local to its domain}

Let \(\cD_{\mathrm{HS}}\subset\cY_{\mathrm{HS}}\) be a gauge-fixed domain on
which a lower evolution \(S_\tau\) is well posed.  Depending on the theorem
being imported, \(S_\tau\) may be a local semiflow rather than a global
semigroup.  We require only
\[
S_0=\operatorname{id},\qquad
S_{\tau+\sigma}=S_\tau S_\sigma
\]
whenever both sides are defined.

\begin{assumption}[Stationary-point identification]
\label{ass:stationary}
On a declared invariant domain \(\cD_0\subset\cD_{\mathrm{HS}}\),
\[
\Fix(S)=\{y\in\cD_0:\cR_{\mathrm{HS}}(y)=0\}.
\]
\end{assumption}

This assumption packages the bundle and global rows that a metric-only
Anomaly flow does not automatically enforce.  In a Fu--Yau ansatz with fixed
valid HYM data, it can be discharged by the corresponding existence and flow
theorems.  It is not currently discharged on the selected physical \(q=79\)
tuple.

\subsection{Why the old twisted complex cannot be used}

\begin{lemma}[Anomaly obstruction to the naive twisted differential]
\label{lem:twisted}
Let \(H\) be a real three-form and define
\(d_H=\dd+H\wedge\) on differential forms.  Then
\[
d_H^2=(\dd H)\wedge.
\]
Consequently \(d_H\) is not a cochain differential on a generic heterotic
background with nonzero Green--Schwarz four-form.
\end{lemma}

\begin{proof}
The graded Leibniz rule gives
\[
(\dd+H\wedge)^2
=\dd H\wedge+H\wedge H\wedge.
\]
Because \(H\) has odd degree, \(H\wedge H=0\).  The stated identity follows.
\end{proof}

Thus the earlier ``twisted harmonic projector'' cannot be justified by
twisted de Rham cohomology unless \(\dd H=0\).  A corrected operator must
instead come from a specified Bismut/Hull covariant Laplacian, an elliptic
deformation complex, or the gauge-fixed Hessian of a selected action.  Its
domain, kernel removal, and spectral gap must be proved for that operator.

\section{The typed MTT bridge}
\label{sec:bridge}

\subsection{Upper data}

Let \(\cU_{\mathrm{MTT}}\) denote an MTT configuration space after quotienting
or fixing its declared redundancies.  A point \(u\) may contain the common
circle line, local \(1<2<3\) filtration, global \(q=79\) carrier data,
projectors, connections, and upper fields.  Let \(R_\tau\) be a selected MTT
stabilization flow on a domain \(\cD_{\mathrm{MTT}}\).

This notation does not assume that such a physical continuum flow has already
been constructed.  It identifies exactly where that future result enters.

\begin{definition}[Typed bridge]
\label{def:bridge}
A typed MTT-to-\HS{} bridge is a map
\[
\mathfrak B:\cD_{\mathrm{MTT}}\longrightarrow\cD_0
\]
whose output is a complete lower datum in the sense of
Definition~\ref{def:lower}.  It must provide the following rows without changing
carrier or source:
\begin{center}
\small
\begin{tabular}{@{}p{0.20\textwidth}p{0.69\textwidth}@{}}
\toprule
Row & Required image under \(\mathfrak B\)\\
\midrule
Carrier & one complex threefold \(X\), complex structure, and \(\Omega\)\\
Metric & positive \(\omega\) and dilaton \(\Phi\)\\
Visible sector & rank-three physical holomorphic bundle and connection\\
Hidden sector & locally free projective rank-nine object, twist, and connection\\
Tangent sector & one declared connection \(\nabla\) and curvature\\
Flux & \(H\), trace convention, gerbe patching, and quantization\\
Dynamics & tangent map carrying the upper vector field to the lower one\\
\bottomrule
\end{tabular}
\end{center}
\end{definition}

The shared circle can enter this bridge as common line-bundle phase or
holonomy data, counted once.  That role does not identify it with Lorentzian
time and does not by itself construct \(X\), \(V_{\mathrm{vis}}\), or the
Green--Schwarz class.

\subsection{One source, two presentations}
\label{sec:source-presentations}

The frozen source-classification and embedding analyses
\cite{FRG1,FRG1E} distinguish two presentations of this program: an augmented
physical bundle deformation germ, and a derived-cohesive germ with a
physical realization map.  This is a classification relative to the stated
routes, not a theorem excluding every conceivable source.  In either
presentation one must preserve provenance, augmentation, nonlinear
operations, metric pairing, gauge action, and all downstream maps.  An
equivalence of tangent dimensions alone preserves none of these six rows.

The elementary embedding \(E\mapsto(E[0],\bar\partial_E)\) preserves the
endomorphism differential graded algebra, its Maurer--Cartan equation, and
gauge transformations.  It preserves adjoints and Hodge operators only after
the same Hermitian structure is carried along.  Conversely, a complex becomes
a physical bundle only after its cohomology is concentrated in degree zero,
is locally free with the required ranks and topology, and carries the correct
twist, augmentation, residual, and readout maps.  For example, the benchmark
\(\kappa_{\mathrm{hol}}\) has \(c_1=-H,c_2=3u\), whereas the physical visible
target has \(c_1=0,c_2=9u,c_3=\pm6\).  Here \(H\) is a divisor class and
\(u\) the normalized base four-class, not the torsion three-form.  The
benchmark is a valid structural source, not that visible bundle.  The
structural repair and Hodge identities are owned by the Cohesive paper
\cite{NeroCohesive2026}; they are not reproved here.

The anti-splicing audit \cite{FRG3A} makes the practical consequence precise.
Metric, bundle, and operator certificates carrying different source-orbit
identifiers cannot be joined without a structure-preserving transport bound
to those sources.  The four candidates examined in that snapshot supplied
no complete physical source.  That is a result about those candidates at
that stage, not a nonexistence theorem for the endpoint.  In particular, its
old missing-hidden-bundle row is superseded by the construction in
Section~\ref{sec:hidden-hym}.

\subsection{Exact intertwining}

\begin{definition}[Flow intertwiner]
\label{def:intertwiner}
The bridge \(\mathfrak B\) intertwines the flows on a common interval \(I\)
containing \([0,\epsilon)\), for some \(\epsilon>0\), if
\begin{equation}
\mathfrak B\circ R_\tau=S_\tau\circ\mathfrak B,
\qquad \tau\in I,
\label{eq:intertwine}
\end{equation}
where both sides are defined.
\end{definition}

Infinitesimally, if the two flows have differentiable vector fields
\(\cF_{\mathrm{MTT}}\) and \(\cF_{\mathrm{HS}}\), exact intertwining requires
\begin{equation}
D\mathfrak B_u\,\cF_{\mathrm{MTT}}(u)
=\cF_{\mathrm{HS}}(\mathfrak B(u)).
\label{eq:generator-intertwine}
\end{equation}
Equation~\eqref{eq:generator-intertwine} is the continuum operator source
obligation.  Matching only fixed-point labels or dimensions does not prove it.

\section{Fixed-point descent}
\label{sec:descent}

\begin{theorem}[Conditional fixed-point descent]
\label{thm:descent}
Let \(R_\tau\), \(S_\tau\), and \(\mathfrak B\) be as above.  Suppose
\(\mathfrak B\) satisfies \cref{eq:intertwine} and
\(u_\ast\in\cD_{\mathrm{MTT}}\) is fixed by \(R_\tau\) for every
\(\tau\in I\).  Then \(\mathfrak B(u_\ast)\) is fixed by \(S_\tau\) for every
\(\tau\in I\).  If Assumption~\ref{ass:stationary} holds, then
\[
\cR_{\mathrm{HS}}(\mathfrak B(u_\ast))=0.
\]
\end{theorem}

\begin{proof}
For every permitted \(\tau\),
\[
S_\tau(\mathfrak B(u_\ast))
=\mathfrak B(R_\tau u_\ast)
=\mathfrak B(u_\ast).
\]
The residual conclusion is exactly Assumption~\ref{ass:stationary}.
Here being fixed for a nontrivial interval from zero means stationary for the
well-posed autonomous local semiflow, and hence fixed throughout its
continuation.  Being fixed at just one isolated time would not suffice.
\end{proof}

\begin{remark}[One-way character]
The theorem is descent, not equivalence.  A lower solution \(y_\ast\) lifts to
an MTT fixed point only if \(y_\ast\in\Ran(\mathfrak B)\) and an upper
preimage is fixed.  Uniqueness of the lift additionally requires control of
the fibers of \(\mathfrak B\).
\end{remark}

\subsection{Residual descent}

Exact commuting diagrams are demanding.  A numerical or perturbative bridge
usually supplies a defect.

\begin{proposition}[Generator-defect bound]
\label{prop:defect}
Suppose \(\mathfrak B\) is differentiable and, on a domain \(\cV\),
\[
\left\|
D\mathfrak B_u\,\cF_{\mathrm{MTT}}(u)
-\cF_{\mathrm{HS}}(\mathfrak B(u))
\right\|_{\cY}
\le \varepsilon .
\]
If \(\cF_{\mathrm{MTT}}(u_\ast)=0\), then
\[
\left\|\cF_{\mathrm{HS}}(\mathfrak B(u_\ast))\right\|_{\cY}
\le\varepsilon .
\]
\end{proposition}

\begin{proof}
Insert \(\cF_{\mathrm{MTT}}(u_\ast)=0\) into the defect inequality.
\end{proof}

This proposition is deliberately modest.  In particular, stationary-point
identification equates zero sets, not norms.  If
\(\cF_{\mathrm{HS}}=-P_y\cR_{\mathrm{HS}}\), an equation estimate requires,
for example,
\[
\|r\|_{\mathrm{res}}\le C_P\|P_y r\|_{\cY}
\quad\text{on the admitted residual range};\qquad
\|\cR_{\mathrm{HS}}(\mathfrak B(u_*))\|_{\mathrm{res}}
\le C_P\varepsilon.
\]
A metric flow that omits bundle rows has no such estimate for those rows.
A small flow residual is not an exact background.  To infer a nearby exact
solution one needs a separate
inverse-function, Newton--Kantorovich, or a posteriori theorem with a
gauge-fixed derivative, inverse bound, nonlinear remainder estimate, and a
verified radius.

\subsection{The full repair operator and the reducing-map requirement}
\label{sec:full-repair}

The residual-Hodge analysis \cite{FRG2} starts at a \emph{zero} of the full
defect.  Write its linearization as \((J,K)\), with \(J\) the cohesive
deformation row and \(K\) the additional geometric rows.  For an orthogonal
unit residual metric, the Hessian of the positive repair cost is
\begin{equation}
 H_{\mathrm{full}}=J^\dagger J+K^\dagger K
 =\Delta_{\cY}+K^\dagger K,\qquad
 \ker H_{\mathrm{full}}=\ker\Delta_{\cY}\cap\ker K.
 \label{eq:full-repair}
\end{equation}
Thus adding equations can remove zero modes.  A bare cohesive compression is
not automatically this full Hessian: the benchmark rank-102 compression
already contains a \(\tfrac14 A_0A_0^\dagger\) contribution beyond its
displayed \(\Delta_Q\) block.  At a nonzero-defect critical point the
residual-weighted second derivative must also be included.

The physical-residual transfer packet \cite{FRG2P} lists six additional rows:
tangent, visible, and hidden moment maps; conformal balance; anomaly/Bianchi;
and \(SU(3)\) normalization.  Together with the base deformation row these
give the seven-row repair problem, after its holomorphic and topological
domain is fixed.  Their cross-couplings matter.  With residual metric
\(W=\left(\begin{smallmatrix}W_0&C\\C^\dagger&W_R\end{smallmatrix}\right)\),
the zero-defect Hessian is instead
\[
J^\dagger W_0J+J^\dagger CK+K^\dagger C^\dagger J+K^\dagger W_RK.
\]
The unit, orthogonal choice in the structural lane is a declared repair
metric, not a derivation of the physical action metric.  Likewise the
rank-102 partition \(3+8+8+80+3\) permits all 25 ordered blocks of a full
matrix, or \(102^2=10404\) positions.  The earlier 19-block mask omitted
2688 permitted positions.  Restoring permission is not computing any of
those matrix entries.

Even an exact residual square does not necessarily transport its Gram
operator.  Let \(T\) and \(S\) be isometric maps on configurations and
residuals, with \(SK_c=K_fT\).  The small example in the same packet is
\[
K_c=[1],\quad K_f=[1\ \tfrac12],\quad
T=\begin{pmatrix}1\\0\end{pmatrix},\quad S=[1].
\]
It gives \(K_f^\dagger K_fT=(1,\tfrac12)^\mathsf{T}\), but
\(TK_c^\dagger K_c=(1,0)^\mathsf{T}\).  Thus the pulled-back cost is exact
while the fine gradient leaves \(\Ran T\).  A reducing or adjoint-compatible
condition, such as \(K_f^\dagger S=TK_c^\dagger\), closes this gap.  For
bounded operators the packet separates residual error and adjoint leakage:
\[
\varepsilon_H\le\varepsilon_0+
(\|K_f\|+\|K_c\|)\varepsilon_K+
\varepsilon_\perp\|K_c\|.
\]
For unbounded continuum operators, a common core and the corresponding
domain estimates are additional requirements.  Cost pullback alone is not
the connection-preserving intertwiner of this paper.

\subsection{What a compiled connection actually transports}
\label{sec:consumer-compiler}

The quantum-mechanics paper owns the projective-module naturality and
Cech compiler used at this interface, in its section \emph{From a geometric
source to a retained operator} \cite{NeroQM2026,ConsumerNaturality,ConsumerCech}.
Given unitary transition functions, a Hermitian metric and a connection,
a partition with \(\sum_i\chi_i^2=1\) constructs an isometric module
embedding \(U\) and the smooth finite-matrix projector \(p=UU^*\).
Its section space is still infinite-dimensional. To preserve the supplied
connection, rather than substitute the Grassmann connection, one uses
\[
 A_0=U^*\dd U,\qquad
 \Gamma=U(A-A_0)U^*,\qquad D_p=p\,\dd+\Gamma .
\]
This is an exact compiler after the geometric inputs are supplied. It
does not choose a physical metric or connection. The external modal flag
also cannot be installed inside the scalar commutant of an irreducible
stable HYM bundle. Cancellation of the hidden central twist in the
rank-80 adjoint explains the ordinary rank-102 deformation complex, not
an ordinary hidden fundamental bundle.

Changing the presentation must transport the intrinsic retained subspace
as well. A fixed Fourier window need not survive a winding change of
frame. For a self-adjoint operator \(H\), exact restriction to \(P\)
requires \((1-P)HP=0\), with compatible domains; otherwise the complementary
resolvent enters through the Feshbach correction. Neither this compiler
nor the abstract Fixed Points analytic spine supplies that reducing
condition for the selected physical operator.

\subsection{A nonlinear strain symbol is not a linear root stack}
\label{sec:consumer-strain}

The same owner supplies the relevant replacement for the excluded linear
root-stack route \cite{NeroQM2026,ConsumerStrain}. The six real root
directions have a different permutation character from two copies of the
three-dimensional permutation representation. They therefore cannot be
identified by an equivariant linear isomorphism. For a Hermitian
three-by-three symbol, the intrinsic strain coordinates instead retain its
three diagonal entries and the three squared edge magnitudes. The complete
relative-phase quotient also retains the triangle phase: its generic
dimension is seven, whereas strain forgets that phase and has six
coordinates. A central shared line does not remove a relative triangle
phase.

On the regular stratum, let \(J\) be the derivative of this nonlinear map
and \(G_Q\) the positive reduced Green operator on the declared source
domain. Minimizing the source quadratic cost at fixed strain gives
\[
 H_{\rm strain}=(JG_QJ^*)^{-1}
\]
on the image, not the restriction of the original Hessian to a guessed
six-dimensional subspace. Any further projection must be applied to the
covariance before inversion. The later support-stratified refinement
settles the zero-edge boundary as well \cite{ConsumerSupportStrain}:
when an edge vanishes, its first derivative vanishes and its squared
magnitude appears at second order. For example, \(r=t^2\) and a source cost
\(t^2/2\) give the normal cost \(r/2\), not a quadratic cost in \(r\).
Tangential shorting remains valid on each support stratum; an ordinary full
six-dimensional boundary Hessian is excluded. A rank-six or the stated
local quarter-turn contract forces the regular stratum. These exact
implications do not supply the physical background, reduced Green operator
or same-source symmetry lift.

\subsection{Finite transfer and the physical source contract}
\label{sec:consumer-transfer}

The Cohesive owner explains why a twisted perfect complex has an ordinary
endomorphism differential graded algebra: the central twist cancels in
endomorphisms and the commutator with its integrable superconnection
squares to zero \cite{NeroCohesive2026,ConsumerCohesive}. On a compact
boundaryless base with chosen Hermitian data, the associated elliptic
Hodge operator has compact resolvent. This global statement does not make
a fiberwise reduced Green operator bounded across a rank-jump locus.
Its Maurer--Cartan residual and gauge row give a positive repair cost
whose exact-zero tangent is the Hodge Laplacian
\cite{ConsumerMC}. Formal Ext/Yoneda transport does not by itself
transport adjoints or identify this cost with the physical signed action;
the reducing-map qualification in \cref{sec:full-repair} remains essential.

The same owner retains the exact finite response contraction and its
nonzero transferred \(m_3,m_4\); it cannot be truncated after the cubic
operation \cite{ConsumerMThree,ConsumerMFour}. All-arity naturality follows
from a product-preserving source map that commutes with the inclusion,
projection and contracting homotopy: it can be moved through every
decorated transfer tree \cite{ConsumerAllArity}. This is an exact
source-map theorem, not an extrapolation from a few tested arities or a
construction of continuum interaction vertices.

Finally, the endpoint-factorization result organizes seven physical
acceptance rows through three structured sources: geometry plus signed
action, spectral synthesis with compatible contraction, and BV-compatible
four-dimensional compactification \cite{ConsumerEndpoint}. Shared-line
functional calculus and Galerkin/Feshbach execution are derived arrows
when those sources agree. Three source packages are not three numerical
parameters, and none is selected by this dependency organization.

\subsection{Scale freedom is not a chosen physical metric}
\label{sec:metric-scale}

The scale-orbit calculation \cite{FRG3B} keeps the source fixed while
multiplying its repair cost by \(\lambda>0\).  The Hessian and its positive
eigenvalues scale by \(\lambda\); the Green operator on the kernel complement
scales by \(\lambda^{-1}\); heat evolution is reparametrized in time.
The kernel and scale-free spectral ratios survive.  Consequently those
ratios cannot select an absolute energy or time normalization.  Independent
relative weights of residual rows are different: they can change the shape
of the operator, not merely its clock.

The metric-commutant packet \cite{FRG3C} illustrates how symmetry can restrict
a \emph{reference} choice.  Sign changes on \(\mathbb R^3\) kill the
off-diagonal entries of an invariant symmetric metric, and permutations make
the three diagonal entries equal.  The invariant positive metrics are
therefore \(\lambda I\).  This finite demonstration does not establish that
the selected physical residual space has this action, or that its metric is
the same reference metric.  Neither calculation turns a positive repair
square into a signed physical action.

\section{What selection would additionally require}
\label{sec:selection}

\begin{proposition}[Contraction-basin criterion]
\label{prop:contraction}
Let \((\cB,d)\) be a nonempty complete invariant subset of
\(\cD_{\mathrm{HS}}\).  If for some \(\tau_0>0\)
\[
d(S_{\tau_0}y,S_{\tau_0}z)\le q\,d(y,z),
\qquad y,z\in\cB,\qquad 0\le q<1,
\]
then \(S_{\tau_0}\) has exactly one fixed point in \(\cB\), and its Picard
iterates converge geometrically to that point.
\end{proposition}

\begin{proof}
This is the Banach fixed-point theorem on \(\cB\).
\end{proof}

The content lies in proving the invariant complete basin and the number
\(q<1\).  Parabolic smoothing alone does not do this.  A schematic estimate
of the form \(C_\Pi M(\tau)e^{L\tau}\) proves contraction only after one
establishes that it is strictly below one for a specified \(\tau\) on a
specified domain.
If \(\cB\) is invariant for the whole semiflow, the unique fixed point of
\(S_{\tau_0}\) is stationary: commutation makes each \(S_s y_*\) another
fixed point in \(\cB\), so uniqueness gives \(S_s y_*=y_*\).  Invariance
under only the sampled time map would not justify this extra conclusion.

\begin{proposition}[Correspondence does not imply selection]
\label{prop:no-selection}
Suppose \(\mathfrak B\) intertwines two flows.  If the lower flow has two fixed
points in \(\mathfrak B(\cD_{\mathrm{MTT}})\), intertwining alone cannot select
one of them.
\end{proposition}

\begin{proof}
The identity map on a space with a flow having two fixed points is already an
intertwiner and selects neither.  Additional basin, initial-data, variational,
or source information is necessary.
\end{proof}

Consequently the phrase ``MTT selects the compactification'' requires more
than \cref{thm:descent}.  It requires a selected upper initial condition or
branch, a complete invariant basin, uniqueness within the physically relevant
quotient, and a proof that no other admissible basin realizes the same
observables.

\subsection{A conditional nontransverse endpoint compiler}
\label{sec:endpoint-compiler}

The nontransverse compiler \cite{FRG3Z} supplies a usable a posteriori contract
for the \emph{full} seven-row residual \(\Psi\).  Work in a fixed gauge slice
on one source, with moduli removed or separately parametrized and the adjoint
obstruction space absent.  At an approximate datum \(s_0\), require
\(A=D\Psi(s_0)\) to be bijective, \(\|A^{-1}\|\le\Gamma\), a derivative
Lipschitz bound \(L>0\), and \(\|\Psi(s_0)\|\le\eta\).  These bounds must
hold on the stated ball of radius \(R\), with a positivity radius
\(r_{\mathrm{pos}}\) and elliptic regularity.  Set
\[
h=\Gamma^2L\eta<\tfrac12,\qquad
r_-={1-\sqrt{1-2h}\over\Gamma L},\qquad
q=\Gamma Lr_-<1.
\]
If \(r_-\le R\) and \(r_-<r_{\mathrm{pos}}\), the frozen-inverse iteration
\(s\mapsto s-A^{-1}\Psi(s)\) contracts the admitted ball and yields a unique
zero there.  The radius equation
\(\Gamma\eta+\tfrac12\Gamma Lr_-^2=r_-\) explains both self-mapping and the
need for a full derivative estimate.  For \(L=0\), the corresponding linear
radius is \(\Gamma\eta\), without division by \(L\).

For scale, the packet's \emph{reference example}, not a \(q=79\)
certificate, uses \(\Gamma=2,L=1/2,\eta=3/16\).  Then
\(h=3/8,r_-=1/2,q=1/2\), fitting inside \(R=1\) and
\(r_{\mathrm{pos}}=3/4\).  This is a bounded arithmetic check of the compiler,
not a computed chamber or connection.  The actual source must supply every
input, preserve its twists and integral classes, and show that gauge
identifications respect the slice.  A zero of this repair residual is an
endpoint under its admitted data; its construction does not identify a
Lorentzian or BV action, choose a global basin, or select the physical source.

\section{Audit of the former selection argument}
\label{sec:audit}

\subsection{The functional was not a selection theorem}

The previous version introduced a functional \(\Xi\) containing the
string-frame curvature term, \(H\)- and Yang--Mills terms, Lagrange
multipliers, and a spectral variance term.  It then claimed:
\[
\operatorname{Crit}(\Xi)
\Longleftrightarrow
\{\text{\HS{} solutions}\},
\qquad
D^2\Xi>0.
\]
Neither implication was established.

First-order supersymmetry equations are not generally identical to the
Euler--Lagrange equations of the ten-dimensional action.  Under additional
instanton and perturbative assumptions, solutions of the supersymmetry and
Bianchi equations can imply the equations of motion, but that is not a
variational equivalence \cite{DelaOssaSvanes2014,AndreasGarciaFernandez2012}.
Second, scalar curvature and multiplier terms do not give an evidently
bounded-below functional.  Third, positivity of a principal elliptic block
does not control all lower-order couplings or moduli.  Finally, a spectral
scalar cannot be asserted to lift every geometric and bundle modulus without
its actual second variation.

Accordingly, this version does not use \(\Xi\), does not claim global
attraction, and does not identify an MTT fixed point with a unique minimizer.
An action-derived repair flow remains a valuable future target, but it must be
constructed before its Hessian or Lyapunov properties are quoted.

\subsection{Connection and torsion bookkeeping}

The gauge-invariant three-form should satisfy both the local Chern--Simons
description and the global differential-cohomological patching law.  Once
\(H\) is used for that gauge-invariant object, the supersymmetry torsion
equation is \(H=\dd^c\omega\) in the chosen convention.  Subtracting the
Chern--Simons terms a second time from the torsion equation double counts
them.

The tangent connection \(\nabla\) also cannot be changed silently between the
torsion equation, Bianchi identity, anomaly flow, and equations-of-motion
claim.  Different choices may define different Hull--Strominger systems
\cite{GarciaFernandez2019}.

\section{Fu--Yau, Iwasawa, and the current MTT evidence}
\label{sec:evidence}

\subsection{What established mathematics supplies}

Fu and Yau constructed solutions on non-K\"ahler torus bundles over K3 under
specific topological, bundle, and analytic hypotheses
\cite{FuYau2008}.  Later work extended and reorganized these constructions,
including solutions with torus symmetry and tangent HYM connections
\cite{FinoGrantcharovVezzoni2021,GarciaFernandez2019}.  The Anomaly flow is
known to converge in a Fu--Yau ansatz for controlled initial data
\cite{PhongPicardZhang2018}.

These results prove that the lower target is mathematically inhabited.  They
do not show that the selected MTT \(q=79\) carrier is one of those solutions,
nor that its upper flow descends to the Anomaly flow.

\subsection{The retained finite and rank-two results}
\label{sec:retained-q79}

Exact arithmetic selects the declared \(q=79\) branch; smooth rank-three
topological candidates have index \(\pm3\); the literal rank-two Cech
witness and finite HYM approximation have a later rank-two Wiener-algebra
existence and local-uniqueness certificate \cite{MTTResultsRepro}.
The analytic certificate is not merely a finite approximation, but it does
not change its bundle rank or carrier.  Similarly, the rank-three topological
candidate does not by itself acquire a holomorphic stable structure.  These
are retained results, not calculations to restart.  The later hidden
rank-nine analysis below changes the evidence table substantially without
closing the physical visible row.

\subsection{Auxiliary spectral data and the first stable source}
\label{sec:aux-spectral}

Let \(S\) be the declared K3 base with classes
\(H^2=2,H\delta=0,\delta^2=-4\), and let \(D_0\) denote the zero section
of its relative elliptic model.  The auxiliary spectral-cutset calculation
\cite{FRSpectral} starts with
\(\operatorname{ch}(\kappa^\sharp)=3-H-3u\).  Its transform has effective
support class \(3H+3D_0\), with smooth finite-flat degree-three members
available by the stated generation and transversality argument.  This is
not the physical visible class \(9H+3D_0\), and it is not a rank-nine hidden
bundle.  The coefficient of \(H\) in a spectral class must not be read as
the bundle rank.  Existence of members also does not choose a particular
member or its twisted line.

The opposite-complex-structure calculation \cite{FRG3E} fixes a genuine sign
issue in the relative transform.  With complex one-forms
\(\eta=\theta-\ii t\) and \(\zeta=s+\ii a\), the Poincare curvature
\[
F_{\mathcal P}=\frac{\ii}{2}
(\bar\eta\wedge\zeta-\eta\wedge\bar\zeta)
\]
has type \((1,1)\); the same-sign choice does not.  The accompanying
\(B\)-shifts are local gerbe curvings, not global ordinary line bundles.
In the declared fiber convention the rank-degree sequence is
\[
(3,0)\longmapsto(3,9)\longmapsto(9,-3)\longmapsto(9,0).
\]
Generically the support splits into ranks and degrees \((1,6)+(2,3)\),
but two exceptional fibers contain a degree-zero piece with shifted
skyscraper transform.  A generic-fiber calculation therefore does not yet
prove local freeness everywhere.

The Hartshorne--Serre source examined in \cite{FRG3F} is an extension
\[
0\longrightarrow\mathcal O^{\oplus2}\longrightarrow F
\longrightarrow\mathcal I_\Gamma(H)\longrightarrow0.
\]
Its two extension images occupy independent character lines.  The locked
divisor-degree restriction excludes positive degree one, while
\(\operatorname{Hom}(F,\mathcal O)=0\) excludes the relevant quotient.
These inputs bound proper rank-one and rank-two slopes by zero, below
\(\mu(F)=2/3\), on the declared balanced ray.  This proves stability of
that source.  Its dual \(\kappa=F^\vee\) is a projective Hermitian--Einstein
\(U(3)\) object with \(c_1=-H\), not the physical visible \(SU(3)\) bundle.

\subsection{What recombination and the qutrit line repair}
\label{sec:recombination}

The two-fiber recombination \cite{FRG3G} replaces a degree-zero summand
\(P=L^{-1}N\) and a degree-three summand \(N\) by the non-split stable
rank-two, degree-three extension represented by the evaluation syzygy
\(Q^\vee\otimes N\).  It introduces no adjustable extension parameter
once the projective class is fixed.  The remaining determinant is a section
of degree two, so its zeros cannot be removed by declaring it an invertible
scalar.  Global twisting and exceptional fibers still require a construction.

The determinantal qutrit packet \cite{FRG3H} supplies a concrete torsion
line.  For \(XZ=\omega ZX\), with \(\omega^3=1\ne\omega\), its pencil
\(M_\omega=uI+vX+wZ\) satisfies
\(\det M_\omega=u^3+v^3+w^3\).  It has corank one on the smooth Fermat
cubic.  The cokernel resolution
\(\mathcal O(-2)^3\to\mathcal O(-1)^3\) identifies a nontrivial
degree-zero order-three line
\(P_\omega=\mathcal O(p_1-p_0)\), where
\(p_0=[-1:0:1]\) and \(p_1=[-\omega:0:1]\), and takes
\(N=L\otimes P_\omega\).  The generator relations explain its finite
projective multiplier.  They do not identify that multiplier with the
analytic gerbe class on a different space.

The forward and reverse Ext computations \cite{FRG3I,FRG3J} illustrate why
orientation of an extension matters.  The forward extension space is
one-dimensional and has a nonzero stable middle term.  The reverse space
has dimension nine, with three copies of each of the three characters.
Nevertheless a reverse extension by a locally free rank-one quotient
splits at the singular stalk into \(\mathcal I\oplus\mathcal O^2\), which
is not locally free along the codimension-two locus.  A larger Ext space
therefore does not imply an acceptable gauge bundle.  This obstruction
excludes the reversed presentation, not the successful forward extension.

\subsection{The derived repairs that do not remove the obstruction}
\label{sec:yoneda}

The derived Yoneda reduction \cite{FRG3K} explains why ordinary scalar and
finite-phase repairs of the locked extension fail.  The required inverse
coefficient line has negative degree and no global section; a Weyl
determinant is instead a positive-degree section and retains rank-drop
zeros.  An odd off-diagonal derived differential using a forward class
\(f\) and reverse class \(\rho\) must also kill its Yoneda products.
The trace relation reduces the two composites to the one relevant
obstruction map.  Finite projective phases cannot bypass that equation.

The monad rank calculation \cite{FRG3L} takes a rank-four quotient of the
nine-dimensional reverse space, leaving a five-dimensional normal space
with character decomposition \(1+2\chi+2\chi^2\).  The clock-fixed
calculation \cite{FRG3M} then isolates a one-dimensional invariant column;
the simultaneous-translation direction maps to zero and is not that
candidate.  The next calculation \cite{FRG3N} evaluates the actual Serre
pairing: the fixed-coordinate normal row is \([1\ 1]\), its boundary is
\((1,-1)\), and the surviving quotient maps nontrivially to the forward
obstruction class.  All three translated columns remain obstructed.  This
closes the proposed locked bidirectional repair, not every derived or
six-by-six presentation of the physical source.

An independent Atiyah/qutrit tensor shortcut \cite{FRG3O} fails at the Chern
sign: after determinant normalization its inverse transform gives
\(c_2=+9u\), rather than the hidden target \(-9u\).  An ambient line does
not fix the locked mixed-class sign.  This is a topological exclusion of
that shortcut, not a negative existence theorem for the hidden bundle.

\subsection{From a perfect complex to a balanced rank-nine bundle}
\label{sec:hidden-object}

The hidden perfect-complex construction \cite{FRG3P} already realizes the
required inverse Chern character \(9+9u\), hence
\((r,c_1,c_2,c_3)=(9,0,-9u,0)\), in a projective derived object involving
a shifted induced term and a zero-section contribution.  It preserves the
categorical and finite multiplier data, but is not a locally free sheaf.
The direct zero-section monad \cite{FRG3Q} cannot repair this: the relevant
fiber \(\operatorname{Hom}\) from a degree-three line to \(\mathcal O\)
vanishes, forcing kernel rank at least three in the proposed map from rank
nine.  Its maximum rank is six.  This excludes that buffer, while preserving
the perfect-complex class.

The global WIT-zero construction \cite{FRG3R} repairs the two exceptional
fibers by stable reduction and determinant-zero-sum clutching weights
\(-1,0,1\).  ``WIT zero'' means that the relative transform is concentrated
in degree zero on \emph{every} fiber; with the verified constant rank and
base-change conditions this gives a locally free projective rank-nine
object on its declared target \(Y=P_{3\delta}\times S^1\).  The analytic
gerbe \(\alpha\), finite qutrit multiplier \(\sigma\), and transported
target twist \(\tau\) have different jobs and cannot be identified merely
because they have order three.  Comparison with a visible bundle on another
presentation of the carrier requires an explicit common-base map.

Local freeness alone did not finish the construction.  The target-slope
audit \cite{FRG3S} finds transformed blocks \((6,+1)+(3,-1)\) before the
support-wide repair.  Their total degree is zero, but the rank-six subobject
has positive slope \(1/6\), excluding that presentation in the declared
adiabatic chamber.  This is not an exclusion for all possible metrics.
The subsequent support-wide balance \cite{FRG3T} uses
\(N\otimes\mathbb C^3_\sigma\): three source factors of rank one and degree
three transform to three \((3,-1)\) factors, then the declared \(B\)-shift
corrects each to \((3,0)\).  Boundary clutching and the Chern class are
preserved.  The resulting holomorphic \(P(3,9)\) object has the
determinant-one reduction specified by that source and is relatively
polystable.  Thus the earlier ``hidden object missing'' headers cannot be
used as its current status.

\subsection{Existential hidden HYM is established}
\label{sec:hidden-hym}

The common-chamber analysis \cite{FRG3D} writes a balanced ray schematically
as
\[
\omega_t=\sqrt t\,\pi^*\omega_H+
t^{-1/2}\ii\zeta\wedge\bar\zeta,\qquad
\deg_t G=tA(G)+B(G).
\]
A strictly negative fiber coefficient and a finite upper slope bound at
\(t=1\) then force negative degree for sufficiently large \(t\).
The comparison is conditional on the actual sheaves, their topology, and
strictness.  For example, a rank bound \(r-1\), negative gap \(a>0\),
and upper slope bound \(M\ge0\) give the sufficient threshold
\(t>1+(r-1)M/a\).  A finite list of factors has a common threshold by taking
the maximum; a visible factor not yet constructed cannot be inserted into
that list.

The direct stable-factor refinement \cite{FRG3U} does not require a reduced
irreducible spectral support.  Each transformed hidden factor has coprime
fiber rank and degree \((3,-1)\), and its corrected proper-subsheaf gap is
\(2/3\).  The packet's provisional numerical slope and chamber slots are
explicit placeholders, not certified values.  The later existential
argument \cite{FRG3V} uses a smooth twisted reference metric and the
compact Chern--Weil projection bound, equivalently the relevant twisted
Harder--Narasimhan upper bound, to obtain a \emph{finite} common \(M\).
For every proper saturated subobject \(G\) of a factor, of rank at most two,
\[
\deg_tG\le2M-\frac23(t-1)<0\qquad(t>1+3M).
\]
The qutrit orbit gives equal total factor degrees, all zero, so one such
\(t\) makes all three factors stable and their rank-nine sum polystable.
The twisted Kobayashi--Hitchin correspondence on the admitted compact
Gauduchon target then supplies a projective HYM connection
\cite{LiYau1987,Perego2019}.

This is an existential HYM result, not an unperformed numerical test.
The absence of a printed numerical \(M\), chamber value, or connection matrix
does not reopen existence.  Equivariant metrics are obtained by transport
around the holomorphic unitary qutrit orbit.  What is still absent is the
same-source visible object, an explicit chamber common to both sectors and
the chosen tangent data, and the coupled Bianchi solution.

\subsection{Two symmetry actions and the anomaly potential}
\label{sec:symmetry-bc}

The hidden adjoint calculation \cite{FRAdjoint} concerns an internal
fiberwise action covering the identity on the base.  Locally
\(W_9=U_3\otimes Q_3\); the projective scalar cancels in
\(\operatorname{End}W_9\).  The honest \(E[3]\) action decomposes the
81-dimensional endomorphism fiber into nine character sectors of rank nine.
Removing the scalar leaves one neutral sector of rank eight and eight
nontrivial sectors of rank nine: \(8+8\cdot9=80\).  Finite Fourier
projectors and an invariant trace pairing are thereby defined.  They do not
compute the physical Hessian on these sectors.

The relative translation result \cite{FRG3W} is different: it moves the
elliptic base point.  Choices of a third root of a translation differ by
three-torsion and line ambiguity, which disappears after projectivization;
the theta-group law supplies coherence.  The finite etale choice torsor
glues over the simply connected K3 base.  This yields a projective action,
not a strict action on vector fibers.  Uniqueness of the stable-factor HYM
connection, with the invariant metric data, makes its quadratic Chern--Weil
form invariant up to gauge.  This does not provide explicit curvature
coefficients.

For the Bianchi cohomology row, the full Bott--Chern argument \cite{FRG3X}
is stronger than an invariant-form computation alone.  On the free
holomorphic principal elliptic action, let \(X_0,Y_0=JX_0\) be the real
generators and \(Z=(X_0-\ii Y_0)/2\).  For a nonzero Fourier character
\((m,n)\), set
\[
\lambda_{mn}=(\ii m+n)/2,\quad
\mu_{mn}=(\ii m-n)/2,\quad
\lambda_{mn}\mu_{mn}=-(m^2+n^2)/4.
\]
Cartan homotopy expresses a Bott--Chern cocycle mode \(\xi_{mn}\) as
\[
\xi_{mn}=(\lambda_{mn}\mu_{mn})^{-1}
\partial\bar\partial\,\iota_{\bar Z}\iota_Z\xi_{mn}.
\]
The inverse quadratic multiplier preserves smooth Fourier summability.
Averaging handles the zero mode using the separately established invariant
injectivity.  Therefore
\(H^{2,2}_{BC}(Y,\mathbb R)\to H^4_{dR}(Y,\mathbb R)\) is injective
on this carrier, not just on its invariant subcomplex.

Once an actual same-source visible--hidden--tangent connection tuple exists,
its real, closed \((2,2)\) anomaly difference with vanishing de Rham class
thus has a \(\dd\dd^c\) potential.  Visible gauge equivariance is no longer
an independent premise for this conclusion, and a separate Aeppli
obstruction is not left over at this linear cohomological stage.  The
potential need not be a positive Hermitian form or satisfy conformal
balance, and it supplies no integral differential trivialization.  Those
nonlinear and integral requirements remain in the endpoint problem.

\subsection{Why the physical endpoint must be nontransverse}
\label{sec:transverse-nogo}

The transverse Fu--Yau exclusion \cite{FRG3Y} concerns the stronger
condition \(F\wedge\omega_B=0\), where \(\omega_B\) is the horizontal
K3 form.  Decompose a two-form into horizontal, mixed, and vertical parts.
On a horizontal space of complex dimension two, wedging a horizontal
one-form with \(\omega_B\) is an isomorphism to horizontal three-forms;
the mixed part is therefore zero.  The vertical part is zero separately.
The remaining curvature is horizontal, so \(F^3=0\).  For the physical
visible \(SU(3)\) bundle this forces \(c_3=0\), contrary to its required
\(c_3=\pm6\).

The excluded condition is not the full six-dimensional HYM equation
\(F\wedge\omega^2=0\).  Established pullback Fu--Yau solutions remain
valid lower-space examples, but cannot be substituted for this visible
bundle.  The nontransverse compiler of
Section~\ref{sec:endpoint-compiler} keeps the mixed and vertical unknowns
and all equations; it is the appropriate conditional analytic contract.

\subsection{Shared-line and integral typing}
\label{sec:shared-line}

The shared-line theta calculation \cite{FRG3AA} works with
\(C_{64}\times E[3]\), of order 576.  Since \(\gcd(64,3)=1\), the mixed
tensor term vanishes; the projective class comes from the \(E[3]\) factor.
For \(W_{\mathrm{sh}}=L_{\mathrm{shared}}\otimes W_9\), an ordinary line
cannot cancel the nontrivial multiplier \(\sigma\).  An inverse-multiplier
module has minimum rank three.  The determinant identity is
\[
\det W_{\mathrm{sh}}
=L_{\mathrm{shared}}^9\otimes\det W_9.
\]
Only on the selected order-two sheet, whose \(C_{64}\) image is
\(\{0,32\}\), does \(L^9=L\).  It is false as an unrestricted identity
of \(C_{64}\) characters.  The resulting flat Spin\(^{c}\) determinant
comparison does not identify nonflat HYM connections.

There is also a substantive correction to earlier ``order-three string
class'' shorthand.  The declared ordinary carrier
\(X=P_\delta\times S^1_{\mathrm{shared}}\) has
\(H^3(X,\mathbb Z)\cong\mathbb Z^{42}\), with no order-three torsion.
The theta class lives in \(H^3(BE[3],\mathbb Z)\) and its pullback to the
quotient stack \([X/E[3]]\), restricting trivially to the ordinary atlas.
The ordinary anomaly cocycle has degree four; its differential
trivializations form a degree-three torsor.  These are distinct from the
equivariant theta class.  A vanishing ambient Brauer value likewise does
not imply that the rank-248 primitive \(\eta_9\) normal function vanishes.
Consumers must use the corrected equivariant statement, not ask for an
impossible ordinary nonzero order-three class on \(X\).

\subsection{The evolving visible-source and transport boundary}
\label{sec:source-frontier}

The long source-frontier chronology \cite{FRFrontier} records mathematical
progress, not merely historical status.  Its later graph/Prym incidence
arguments establish available smooth finite-flat supports and formal local
lifts, while excluding a particular fixed-graph shortcut.  They do not yet
choose the twisted line with the required primitive Deligne zero.  The
locked tangent rank \(122=33+89\) in an ambient rank-248 problem leaves
126 normal directions; finite rank checks and a nonzero integral minor can
control a local response without evaluating the nonlinear primitive normal
function.  A zero modular obstruction is not an exact zero.  A nonzero
modular obstruction excludes an exact zero only for the same integral
model, field presentation, and declared specialization.

The integral period analysis in the same source distinguishes the image
\(4\mathbb Z\) detected by algebraic product tubes from the full integral
pairing image \(\mathbb Z\).  Thus a root readout \(4\beta\) vanishing
can leave four-torsion in \(\beta\).  A direct Deligne zero requires
integral period membership, while a nonzero certificate requires a dual
character integral on the full relevant lattice, or a proved saturated
detector.  Rational span alone is insufficient.  In the primitive setting,
the integral ambient quotient is the dual lattice \(V^\vee\), not
automatically the vanishing lattice \(V\); their finite discriminant
quotient must be retained.  The Cohesive discussion \cite{NeroCohesive2026}
also retains the corrected CBF restriction at T69, the local nature of
T70--T73 as distinct from global Picard/BHT identification, and the fact
that an ordinary intermediate-Jacobian Neron model requires its own
weight and admissibility hypotheses.  None of these distinctions is
removed by a shared theta multiplier.

The chronology's exact smooth starting family and discriminant-avoiding
path construction provide geometric inputs for transport, not an enclosed
global period integral.  Its static Schur reduction from 1510 coordinates
to 249 failed the later tangency/Riccati tests as a differential subsystem.
A moving basis would require the connection term \(B^{-1}\dd B/\dd a\).
Matrix-free sector evaluation, finite jet ranks, and a small static solve
therefore cannot stand in for the full transported source.  Later finite
and certified local \(q=79\) results remain established at their current
authority; the still-open target is a useful global path with the joint
primitive readout, not a repetition of those local calculations.

The chronology also retains a source-exact Cox presentation and finite
component calculations while rejecting specific scalar/apolar and
derivative-closure intertwiners.  The 48-place fiber results are exact over
their stated finite fields, not an equality of characteristic-zero
primitive source planes.  The later rank-18 source injection is a
same-field statement, not an identification of every target plane.  In
particular, substituting \(\mathbb F_{11}\) coordinates into an
\(\mathbb F_{21817}\) calculation is not transport.

The separate arithmetic correction \cite{FRG3GM} retracts one concrete
fresh-prime rank jump.  Calls labelled 79903, 81359, and 82963 inherited
the old modulus 33863.  The reported rank/nullity \(4752/1270\) and the
associated component-23/24 changes are therefore not new-prime results.
Corrected replay at the archived \(p=21817\) chart and an actual
\(p=79903\) chart restores the scoped finite rank/nullity \(4748/1274\).
The valid characteristic-zero conclusion remains the bound
\[
4748\le\operatorname{rank}_{\mathbb Q(\gamma)}Q\le5102,
\]
not equality with the finite rank.  Unaffected fixed-minor and separate
finite-component results are retained at their own tiers.  This is a
specific arithmetic/provenance correction, not a blanket downgrade of the
evolving source.

\subsection{Global support and the selected transport certificates}
\label{sec:consumer-global}

The Cohesive owner's Fitting-descent statement supplies ordinary ideal
sheaves for the cohomology of the twisted benchmark: invertible transition
lines change local minor generators by units, not their ideals
\cite{NeroCohesive2026,ConsumerFitting}. Its leading divisor is
\(3H+3D_0\), agreeing with the auxiliary spectral class in
\cref{sec:aux-spectral}, not the physical \(9H+3D_0\). Generic corank three
does not produce a spectral line, trivialize the gerbe, or execute the
literal connection overlaps. This is useful global support information
without any of those promotions.

The same owner's large-gauge result concerns a different question
\cite{ConsumerBK3}. The selected real K3 period plane has zero
large-gauge kernel; its two-real-parameter image is nevertheless dense
and nonclosed in a twenty-dimensional torus. Rank zero counts
identifications, not remaining parameters. Numerical periods and physical
normalization are not recovered from this injectivity theorem, and it
does not trivialize the unrelated visible obstruction \(\beta_C\).

The flux paper owns the selected eta9 transport calculations
\cite{NeroFlux2026}. The completed original-Jacobian campaign supplies
all 30 groups and 225 support columns as exact polynomial reduction
identities \cite{ConsumerGateOne}. This closes that finite algebraic
task, not the global source path. The three-cycle calculation has
Gram matrix \(-2I_3\): it proves independence, but its affine coefficients
are rational coboundaries on the tested subsystem
\cite{ConsumerThreeCycle}. Thus this complete subsystem is
non-detecting; it neither proves integral coboundary membership nor
trivializes the global affine class.

There is also a certified positive-width local Gauss--Manin source chart,
with its correlated lift, rather than merely a single-fiber calculation
\cite{ConsumerTube}. The consolidation is a hash-bound record of those
certificates, not an independent interval replay. A useful global path,
integral marking and physical readout remain separate. Conversely, the
same-source B89 detector pairs nontrivially with the affine return modulo
two and therefore rejects that candidate from the \(\beta_C=0\) locus
\cite{ConsumerB89}. It does not show that the integral class has exact
order two, or supply a replacement visible bundle. These distinctions
explain how completed local work and hidden HYM existence coexist with
the open common endpoint.

\subsection{Current completion gates}
\label{sec:current-gates}

The current authority distinguishes these rows.  ``Established'' below
means established at the named source tier, not independent verification by
this editorial review.
\begin{center}
\small
\begin{tabular}{@{}p{0.10\textwidth}p{0.55\textwidth}p{0.25\textwidth}@{}}
\toprule
Gate & Object or implication & Current scope\\
\midrule
V1 & selected visible rank-three twisted Prym bundle, index \(\pm3\)
   & open\\
V2 & locally free holomorphic hidden \(P(3,9)\) object
   & established\\
H2 & projective hidden HYM connection on its declared target
   & existentially established\\
H1 & visible, hidden, and tangent instantons on one common metric/base
   & open\\
A0 & exact real anomaly source implies a \(\dd\dd^c\) potential
   & conditional on the source pair\\
A1 & positive conformally balanced nontransverse endpoint
   & open; compiler conditional\\
Q1 & ordinary differential Green--Schwarz trivialization and equivariant lift
   & open\\
B1 & same-source connection/domain-preserving continuum intertwiner
   & open\\
\bottomrule
\end{tabular}
\end{center}
The open \(B.\mathrm{HS}.01\) and \(B.\mathrm{GEO}.01\) obligations are
therefore narrower than ``no hidden bundle or HYM result.''  They cannot be
closed by assembling established rows from different source orbits.

\subsection{The corrected role of Iwasawa and Lens--Nil}

The Iwasawa manifold remains a useful source of explicit balanced
\(SU(3)\)-structure calculations.  The companion bundle audit
\cite{NeroFluxAudit2026} shows, however, that the earlier printed rank-three
bundle, Bianchi match, and selected-background conclusion do not survive.
Other valid Iwasawa Hull--Strominger solutions exist in the literature; their
existence does not repair that specific MTT construction.

The full Qa/SU3 support chronology \cite{FRQa} must be read past its early
failed connection and stale status headings.  Its corrected Maurer--Cartan
representative, and the later finding that the displayed family is a single
complex gauge orbit with a noncentral commutant, do not supply a stable
physical color source.  A branch-specific closed-curvature representative
does not replace the generic \(\dd H\ne0\) anomaly bookkeeping of
Lemma~\ref{lem:twisted}.  Those are exclusions of particular proposals,
not of every Iwasawa calculation.

The later finite internal computation is retained.  For instance its
declared quotient gives the charge Gram matrix
\[
H_{\mathrm{sel}}=
\begin{pmatrix}26&-3&0\\-3&10&0\\0&0&8\end{pmatrix},
\qquad \det H_{\mathrm{sel}}=2008.
\]
This explains its internal logarithmic finite part after the stated zero-mode
quotient.  It is not the smooth physical color Laplacian.  Similarly, a
phase-preserving injection into a 27-dimensional label space need not
intertwine a positive operator or its determinant.  The later oriented
functional driver with normalized value \(N_{\alpha_1}(h_{\mathrm{ext}})=1\)
is a source-scoped tangent result, not a physical action amplitude.  Its
stationary and derivative conclusions require the declared driver and
gauge transport.  The chronology's old ``SM open'' summaries do not
override the current established finite SM, shared-primitive, Yukawa, and
electroweak results.  None of those finite results is reopened here.

Likewise,
\[
\text{local Circle--Lens--Nil filtration}
\not\equiv L(3,1)\times\mathrm{Nil}_3
\not\equiv X_{79}.
\]
Lens and Nil may remain auxiliary finite, spectral, or transport labels.
They are not a substitute for the integrable complex geometry and global
bundle data in Definition~\ref{def:lower}.

\section{Worldsheet and physical interpretation}
\label{sec:physical}

A first-order Hull--Strominger background is an important string-theoretic
object, but it is not by itself a complete four-dimensional model.  A
worldsheet completion requires the relevant anomaly, conformal, modular,
factorization, and GSO data.  Four-dimensional physics additionally requires
the massless cohomology, representation embedding, normalized kinetic terms,
overlap integrals, moduli treatment, and quantum corrections.

The fixed-point bridge therefore has a precise interpretation:
\begin{itemize}[leftmargin=1.5em]
\item it can explain how an upper MTT equilibrium becomes a lower geometric
      equilibrium;
\item it can transport a certified residual and its error budget;
\item it cannot create missing bundles or flux data;
\item it cannot turn a conditional encoding into source selection; and
\item it does not derive low-energy masses or couplings without additional
      normalized overlap and transport maps.
\end{itemize}

\section{A concrete completion program}
\label{sec:frontier}

The revised bridge suggests an efficient order of work.
\begin{enumerate}[leftmargin=1.7em]
\item Supply the selected visible twisted Prym datum, including its primitive
      Deligne-zero condition, and construct the physical rank-three bundle
      with the required topology and stability.
\item Retain the established hidden projective rank-nine object and
      existential HYM connection.  Place the visible, hidden, and tangent
      instantons on one declared common base and prove a common chamber.
\item Use the full Bott--Chern reduction for the actual anomaly source and
      establish a positive conformally balanced nontransverse solution,
      for example by supplying the full compiler's analytic bounds.
\item Construct the ordinary integral differential Green--Schwarz
      trivialization and the correctly typed equivariant theta comparison;
      do not replace either with a flat determinant identity.
\item Define the MTT configuration space and its selected continuum vector
      field on that same tuple.
\item Define every row of \(\mathfrak B\), including connection and domain
      maps.
\item Prove \cref{eq:generator-intertwine}, or emit a certified defect bound.
\item Use \cref{thm:descent} or \cref{prop:defect}; if selection is claimed,
      independently prove the contraction-basin hypotheses of
      \cref{prop:contraction}.
\end{enumerate}

The Anomaly flow supplies an established lower fixed-point language under
its analytic hypotheses.  The full residual compiler gives another
conditional local route once all same-source data and bounds are supplied.
Neither route requires reopening established finite SM, operational quantum
mechanics, hidden HYM, or local \(q=79\) conclusions.  Measurement in any
eventual physical model remains ordinary physics; this geometric completion
contract does not introduce a separate actualization law.

\section{Conclusion}

The viable relationship between MTT and the Hull--Strominger system is a
conditional fixed-point correspondence.  Its central object is not an
assumed selection potential but a typed, connection-preserving flow
intertwiner.  If that intertwiner and one upper fixed point are supplied,
fixed-point descent is exact.  If only a defect estimate is supplied, the
result is an approximate lower solution until an a posteriori theorem closes
the residual.

The source now includes a holomorphic hidden \(P(3,9)\) object, an
existential projective HYM connection, and a full Bott--Chern anomaly-potential
reduction on the declared \(q=79\) target.  These coexist with the retained
finite, topological, and rank-two results.  The physical visible rank-three
bundle and its integral normal-function condition, a common positive
nontransverse endpoint, and the connection-preserving intertwiner remain
distinct obligations.  Preserving those distinctions uses the successful
construction without reviving its excluded predecessors or overstating
physical selection.



\bibliographystyle{plain}

% BEGIN MTT MANAGED COMPUTATIONAL EVIDENCE
\section*{Computational Evidence and Reproducibility}
\label{sec:reproducibility}
The retained arithmetic, literal rank-two Cech, and rank-two Wiener results
remain at their stated tiers.  The new contextual citations bind all 37
assigned frozen sources, including the two evolving chronologies, to the
curated manifest commit below.  Their explained imports in
Sections~\ref{sec:source-presentations}--\ref{sec:source-frontier} do not
duplicate their proof bodies or promote a physical endpoint.

The referenced rows are frozen to the curated results repository state identified below. Hashes are grouped in eight-character blocks for line breaking.
\begin{quote}\small
\textbf{Repository:} \url{https://github.com/PeterNero/mtt-results-repro}\\
\textbf{Commit:} \texttt{f141a20e a23c5c3f f19cc216 1c0e226e 29ade8a7}\\
\textbf{Manifest:} \path{release/result_manifest.json}
\end{quote}

Tier labels are quoted verbatim from that manifest. A row used directly supports only the specific computational statement identified above; a corpus-state cross-check does not prove this paper's local theorems; and an open row is evidence of an unresolved obligation, never of closure.

The upstream hash lock \cite{FRLock} is retained as historical provenance:
it identifies the source inputs available to that snapshot, not an extra
theorem or current authority.  Its entire map was reviewed, not dismissed
from an early status heading.  Current A/B authority and later scoped
refinements govern present conclusions.  By contrast the Qa/SU3 and unified
frontier chronologies contain substantive later results and exclusions;
their contextual imports are not classified wholesale as historical.

The local research-integration record and the source-bound editorial review
fragment give exact artifact paths, SHA-256 hashes, and manuscript anchors
for each assignment.  Full source reading and small checks of the displayed
matrix, slope, and compiler examples are editorial validation, not
independent verification of a boolean packet status.  No scientific source
calculation was rerun in this revision.

\par\begingroup\dimen0=\pagegoal\advance\dimen0 by -\pagetotal\ifdim\dimen0<7\baselineskip\newpage\fi\endgroup
\paragraph{Rows used directly in this paper.}
\begin{itemize}
\setlength{\itemsep}{0.25em}
\item \path{A19/hym_wiener_contraction} (\emph{numeric certified}).\par Weighted-theta Fourier-tail and Wiener contraction certificate.
\item \path{A07/literal_cech_witness} (\emph{derived exact}).\par Literal 81-entry, 729-cocycle finite Cech witness.
\item \path{A11/q79_exact_audit} (\emph{derived exact}).\par Executable q=79 exact-branch audit.
\item \path{A11/q79_exact_theorem} (\emph{derived exact}).\par CRT q=79 theorem on the selected exact branch.
\end{itemize}

No imported row changes theorem ownership or promotes a neighboring claim: all local statements retain their stated hypotheses, domains, and limitations.
% END MTT MANAGED COMPUTATIONAL EVIDENCE

\bibliography{main}

\end{document}
